% يجب الترجمة بمحرك LuaLaTeX  (Menu -> Compiler -> LuaLaTeX)
% اختبار العلوم - الصف الثامن، ويليه الإجابة النموذجية في صفحة جديدة
\documentclass[12pt, a4paper]{article}
\usepackage[a4paper, top=1.2cm, bottom=1.4cm, left=1.4cm, right=1.4cm]{geometry}
\usepackage{amsmath, amssymb}
\usepackage{fontspec}
\usepackage[bidi=basic]{babel}
\babelprovide[main, import, onchar=ids fonts]{arabic}
\babelprovide[import, onchar=ids fonts]{english}
\babelfont{rm}{Noto Serif}
\babelfont[arabic]{rm}{Amiri}
\usepackage[version=4]{mhchem}
\usepackage{enumitem}
\usepackage{tabularx}
\usepackage{array}
\usepackage{graphicx}
\usepackage{tikz}
\usetikzlibrary{arrows.meta}
\usepackage{eso-pic}

\setlength{\parindent}{0pt}
\newcommand{\LTRfig}[1]{\mbox{\foreignlanguage{english}{#1}}}
\newcommand{\en}[1]{\foreignlanguage{english}{#1}}

% عنوان سؤال داخل شريط رمادي مع العلامة على اليسار
\newcommand{\question}[2]{%
  \par\vspace{0.2cm}\noindent
  \colorbox{gray!15}{\parbox{\dimexpr\linewidth-2\fboxsep\relax}{\textbf{\large #1}\hfill\mbox{\textbf{(#2)}}}}%
  \par\vspace{0.12cm}}
% سطر منقط يكمل السطر، وسطر منقط كامل
\newcommand{\fillline}{\leavevmode\xleaders\hbox to 0.35em{\hss.\hss}\hfill\kern0pt}
\newcommand{\dotline}{\par\vspace{0.2cm}\noindent\fillline\par}
% فراغ منقط بطول محدد
\newcommand{\dotsblank}[1][3.2cm]{\makebox[#1]{\xleaders\hbox to 0.35em{\hss.\hss}\hfill}}
% (للإجابة) علامة الفقرة والإجابة داخل إطار
\newcommand{\m}[1]{\hfill\mbox{\textcolor{blue!60!black}{\textbf{[#1]}}}}
\newcommand{\ans}[1]{{\setlength{\fboxsep}{1.5pt}\fbox{\textbf{#1}}}}

% الترويسة (تُستعمل للاختبار وللإجابة النموذجية)
\newcommand{\header}[1]{%
\noindent
\begin{minipage}[t]{0.37\textwidth}
\vspace{0pt}
\textbf{المملكة الأردنية الهاشمية} \\
\textbf{وزارة التربية والتعليم} \\
\textbf{مدرسة جحفية الشاملة للبنات} \\
\textbf{العام الدراسي: 2026 / 2027}
\end{minipage}%
\begin{minipage}[t]{0.26\textwidth}
\vspace{0pt}
\centering
\IfFileExists{logo.png}{\includegraphics[width=2.7cm]{logo.png}}{}
\end{minipage}%
\begin{minipage}[t]{0.37\textwidth}
\vspace{0pt}
\textbf{المبحث: العلوم} \\
\textbf{الصف: الثامن} \\
\textbf{الزمن:} \rule{2.5cm}{0.4pt} \\
\textbf{العلامة الكلية: 80 علامة}
\end{minipage}

\vspace{0.2cm}
\hrule height 1.2pt
\vspace{0.06cm}
\hrule height 0.4pt
\vspace{0.12cm}

\begin{center}
\textbf{\LARGE #1}
\end{center}
\vspace{-0.2cm}}

\begin{document}

% إطار الصفحات
\AddToShipoutPictureBG{%
  \LTRfig{%
  \begin{tikzpicture}[remember picture, overlay]
    \draw[line width=1.2pt] ([shift={(0.7cm,-0.7cm)}]current page.north west) rectangle ([shift={(-0.7cm,0.7cm)}]current page.south east);
  \end{tikzpicture}}%
}

% ####################################################################
% ############################ الاختبار ##############################
% ####################################################################
\header{اختبار مادة العلوم}

% بيانات الطالبة
\noindent\fbox{\parbox{\dimexpr\linewidth-2\fboxsep-2\fboxrule\relax}{%
\vspace{0.15cm}
\textbf{اسم الطالبة:} \rule{6.2cm}{0.4pt} \hfill
\textbf{الشعبة:} \rule{2cm}{0.4pt} \hfill
\textbf{التاريخ:} \rule{0.9cm}{0.4pt} / \rule{0.9cm}{0.4pt} / \LTRfig{202\rule{0.6cm}{0.4pt}}
\vspace{0.1cm}}}

\vspace{0.12cm}
\begin{center}
\renewcommand{\arraystretch}{1.2}
\begin{tabular}{|>{\bfseries}c|*{6}{>{\centering\arraybackslash}p{1.35cm}|}>{\centering\arraybackslash}p{1.5cm}|}
\hline
السؤال & الأول & الثاني & الثالث & الرابع & الخامس & السادس & المجموع \\ \hline
العلامة & 30 & 6 & 24 & 10 & 4 & 6 & \textbf{80} \\ \hline
علامة الطالبة & & & & & & & \\ \hline
\end{tabular}
\end{center}
\vspace{-0.2cm}
\noindent\textbf{تعليمات:} أجيبي عن جميع الأسئلة الآتية.

% ============================ السؤال الأول ============================
\question{السؤال الأول: أكملي الفراغ بما يناسبه}{30 علامة}

\begin{enumerate}[label=\textbf{\arabic*-}, itemsep=0.42cm, leftmargin=*]
\item يُقصد بتلوث المياه: \fillline \dotline
\item دور العمليات الجيولوجية في تشكيل الموارد المعدنية: \par\vspace{0.2cm}
\textbf{1)} \fillline \quad \textbf{2)} \fillline
\item تتحرك الصفائح التكتونية بالنسبة إلى بعضها، وبناءً على ذلك تتكون ثلاثة أنواع من الحدود هي: \par\vspace{0.2cm}
\textbf{1)} \fillline \quad \textbf{2)} \fillline \quad \textbf{3)} \fillline
\item تتكون النيوكليوتيدات من: \par\vspace{0.2cm}
\textbf{1)} \fillline \quad \textbf{2)} \fillline \quad \textbf{3)} \fillline
\item ترتبط القاعدة \en{(G)} مع القاعدة \en{(C)} بـ \dotsblank\ روابط هيدروجينية.
\item يحدث الانقسام المتساوي في الخلايا \fillline
\item بينما يحدث الانقسام المنصف في الخلايا \fillline
\item يمر الانقسام المتساوي بأطوار عدة، منها: \par\vspace{0.2cm}
\textbf{1)} \fillline \quad \textbf{2)} \fillline \quad \textbf{3)} \fillline
\end{enumerate}

\newpage
% ============================ السؤال الثاني ============================
\question{السؤال الثاني: الوراثة}{6 علامات}

لقّح مندل نباتَي بازيلاء، أحدهما طويل الساق غير متماثل الأليلات، والآخر طويل الساق غير متماثل الأليلات.
اكتبي الطرز الجينية والشكلية للأبناء باستخدام مربع بانيت، علمًا بأن صفة طول الساق سائدة.

\vspace{0.2cm}
\noindent
\begin{minipage}[t]{0.5\linewidth}
\vspace{0pt}
رمز الأليل السائد: \dotsblank[2cm] \qquad رمز الأليل المتنحي: \dotsblank[2cm]
\par\vspace{0.45cm}
الطرز الجينية للأبوين: \dotsblank[4.5cm]
\par\vspace{0.45cm}
الطرز الجينية للأبناء: \fillline
\par\vspace{0.45cm}
الطرز الشكلية للأبناء: \fillline
\par\vspace{0.45cm}
النسبة: \fillline
\end{minipage}%
\hfill
\begin{minipage}[t]{0.4\linewidth}
\vspace{0pt}
\centering
\LTRfig{%
\begin{tikzpicture}
  \draw[thick] (0,0) grid[step=1.4] (4.2,-4.2);
  \fill[gray!15] (0,0) rectangle (1.4,-1.4);
  \draw[thick] (0,0) grid[step=1.4] (4.2,-4.2);
  \node[font=\small] at (2.8,0.35) {\foreignlanguage{arabic}{جاميتات الأب}};
  \node[font=\small, rotate=90] at (-0.35,-2.8) {\foreignlanguage{arabic}{جاميتات الأم}};
\end{tikzpicture}}
\end{minipage}

% ============================ السؤال الثالث ============================
\question{السؤال الثالث: لديك العناصر الآتية}{24 علامة}

\begin{center}
\LTRfig{\Large\textbf{\ce{_{32}Ge} \qquad\quad \ce{_{19}K} \qquad\quad \ce{_{7}N}}}
\end{center}
\vspace{-0.15cm}

\textbf{(أ)} حددي لكل عنصر ما يلي، بإكمال الجدول:
\par\vspace{0.15cm}
\begin{center}
\renewcommand{\arraystretch}{1.9}
\begin{tabular}{|c|>{\centering\arraybackslash}p{1.9cm}|>{\centering\arraybackslash}p{1.9cm}|>{\centering\arraybackslash}p{2.4cm}|>{\centering\arraybackslash}p{2.4cm}|>{\centering\arraybackslash}p{2.4cm}|}
\hline
\textbf{العنصر} & \textbf{1- الدورة} & \textbf{2- المجموعة} & \textbf{3- عدد إلكترونات التكافؤ} & \textbf{4- عدد مستويات الطاقة} & \textbf{5- الأيون المتوقع} \\ \hline
\en{\ce{_{32}Ge}} & & & & & \\ \hline
\en{\ce{_{19}K}} & & & & & \\ \hline
\en{\ce{_{7}N}} & & & & & \\ \hline
\end{tabular}
\end{center}

\textbf{6-} اسم المجموعة التي تكتمل فيها مستويات الطاقة الخارجية: \fillline
\par\vspace{0.3cm}
\textbf{7-} أي هذه العناصر ينتمي إلى الفلزات؟ \fillline
\par\vspace{0.35cm}
\textbf{(ب)} حددي عدد كل من: النيوترونات \en{$n^{0}$} ، والبروتونات \en{$p^{+}$} ، والإلكترونات \en{$e^{-}$} لذرة عددها الكتلي \en{(32)} وعددها الذري \en{(15)}.
\par\vspace{0.35cm}
\en{$p^{+}$} = \dotsblank \qquad\quad \en{$e^{-}$} = \dotsblank \qquad\quad \en{$n^{0}$} = \dotsblank

\newpage
% ============================ السؤال الرابع ============================
\question{السؤال الرابع: وضّحي المقصود بكل مما يأتي}{10 علامات}

\begin{enumerate}[label=\textbf{\arabic*)}, itemsep=0.2cm, leftmargin=*]
\item الضغط: \fillline \dotline
\item قاعدة باسكال: \fillline \dotline
\item مبدأ برنولي: \fillline \dotline
\item قوة الطفو: \fillline \dotline
\item قاعدة أرخميدس: \fillline \dotline
\end{enumerate}

% ============================ السؤال الخامس ============================
\question{السؤال الخامس: الرافعة الهيدروليكية}{4 علامات}

في رافعة هيدروليكية، كانت مساحة سطح المكبس الصغير \en{$0.4\ \mathrm{m^2}$}، ومساحة سطح المكبس الكبير \en{$1.6\ \mathrm{m^2}$}.
فما مقدار القوة اللازم التأثير بها في المكبس الصغير لرفع سيارة وزنها \en{$24000\ \mathrm{N}$}؟
\vspace{3.8cm}

% ============================ السؤال السادس ============================
\question{السؤال السادس: اشرحي بالرسم العلاقة بين قوة الطفو والوزن}{6 علامات}

\vspace{0.1cm}
\noindent\fbox{\parbox[t][6.3cm][t]{\dimexpr\linewidth-2\fboxsep-2\fboxrule\relax}{\small مساحة الرسم:}}

\vfill
\noindent\textbf{\large مع دعائي لكنّ بالنجاح} \hfill \textbf{\large المعلمة: دينا الشرمان}

% ####################################################################
% ######################## الإجابة النموذجية #########################
% ####################################################################
\newpage
\header{الإجابة النموذجية لاختبار العلوم}
\noindent{\small الأرقام بين الأقواس المربعة \textcolor{blue!60!black}{\textbf{[ ]}} هي علامة كل فقرة. تُقبل الإجابات الصحيحة الأخرى المكافئة.}

% ============================ السؤال الأول ============================
\question{السؤال الأول: أكملي الفراغ بما يناسبه}{30 علامة: علامتان لكل فراغ}

\begin{enumerate}[label=\textbf{\arabic*-}, itemsep=0.08cm, leftmargin=*]
\item تلوث المياه: \ans{أي تغيّر في خصائص الماء الفيزيائية أو الكيميائية أو الحيوية يجعله ضارًّا بالكائنات الحية أو غير صالح للاستخدام.} \m{2}
\item دور العمليات الجيولوجية في تشكيل الموارد المعدنية: \m{4}
\par \textbf{1)} \ans{تبلور الصهارة وتبرّدها (العمليات النارية) فتتكوّن المعادن وتتركّز.}
\par \textbf{2)} \ans{التجوية والتعرية والترسيب (العمليات الرسوبية) فتتجمّع المعادن في رواسب.}
\item أنواع حدود الصفائح: \ans{حدود متباعدة} ، \ans{حدود متقاربة} ، \ans{حدود تحويلية (انزلاقية)} \m{6}
\item تتكون النيوكليوتيدات من: \ans{مجموعة فوسفات} ، \ans{سكر خماسي (رايبوز منقوص الأكسجين)} ، \ans{قاعدة نيتروجينية} \m{6}
\item ترتبط القاعدة \en{(G)} مع \en{(C)} بـ \ans{ثلاث} روابط هيدروجينية. \m{2}
\item يحدث الانقسام المتساوي في الخلايا \ans{الجسمية} \m{2}
\item يحدث الانقسام المنصف في الخلايا \ans{الجنسية (التي تكوّن الجاميتات في المناسل)} \m{2}
\item من أطوار الانقسام المتساوي (أي ثلاثة): \ans{التمهيدي} ، \ans{الاستوائي} ، \ans{الانفصالي} ، \ans{النهائي} \m{6}
\end{enumerate}

% ============================ السؤال الثاني ============================
\question{السؤال الثاني: الوراثة}{6 علامات}

\noindent
\begin{minipage}[t]{0.58\linewidth}
\vspace{0pt}
الأليل السائد (طويل): \ans{\en{T}} \quad الأليل المتنحي (قصير): \ans{\en{t}}
\par\vspace{0.12cm}
الطرز الجينية للأبوين: \ans{\en{Tt $\times$ Tt}} \m{1}
\par\vspace{0.12cm}
الجاميتات: \en{T} ، \en{t} من كل أب \m{1}
\par\vspace{0.12cm}
الطرز الجينية للأبناء: \ans{\en{TT}} \ ، \ \ans{\en{Tt}} \ ، \ \ans{\en{tt}} \ بنسبة \ans{\en{1 : 2 : 1}} \m{2}
\par\vspace{0.12cm}
الطرز الشكلية: \ans{طويل الساق : قصير الساق} بنسبة \ans{\en{3 : 1}} \m{2}
\end{minipage}%
\hfill
\begin{minipage}[t]{0.35\linewidth}
\vspace{0pt}
\centering
\LTRfig{%
\begin{tikzpicture}[scale=0.85]
  \fill[gray!15] (0,0) rectangle (1.4,-1.4);
  \draw[thick] (0,0) grid[step=1.4] (4.2,-4.2);
  \node[font=\Large\bfseries] at (2.1,-0.7) {T};
  \node[font=\Large\bfseries] at (3.5,-0.7) {t};
  \node[font=\Large\bfseries] at (0.7,-2.1) {T};
  \node[font=\Large\bfseries] at (0.7,-3.5) {t};
  \node[font=\large] at (2.1,-2.1) {TT};
  \node[font=\large] at (3.5,-2.1) {Tt};
  \node[font=\large] at (2.1,-3.5) {Tt};
  \node[font=\large] at (3.5,-3.5) {tt};
\end{tikzpicture}}
\end{minipage}

% ============================ السؤال الثالث ============================
\question{السؤال الثالث: العناصر}{24 علامة}

\textbf{(أ)} {\small التوزيع الإلكتروني: \en{\ce{Ge}}: 2، 8، 18، 4 \quad \en{\ce{K}}: 2، 8، 8، 1 \quad \en{\ce{N}}: 2، 5} \m{15: علامة لكل خانة}
\begin{center}
\vspace{-0.2cm}
\renewcommand{\arraystretch}{1.45}
\begin{tabular}{|c|>{\centering\arraybackslash}p{1.9cm}|>{\centering\arraybackslash}p{2.1cm}|>{\centering\arraybackslash}p{2.3cm}|>{\centering\arraybackslash}p{2.3cm}|>{\centering\arraybackslash}p{2.4cm}|}
\hline
\textbf{العنصر} & \textbf{الدورة} & \textbf{المجموعة} & \textbf{إلكترونات التكافؤ} & \textbf{مستويات الطاقة} & \textbf{الأيون المتوقع} \\ \hline
\en{\ce{_{32}Ge}} & الرابعة & \en{4A} (14) & 4 & 4 & \small لا يكوّن أيونًا بسهولة* \\ \hline
\en{\ce{_{19}K}} & الرابعة & \en{1A} (1) & 1 & 4 & \en{\ce{K+}} \\ \hline
\en{\ce{_{7}N}} & الثانية & \en{5A} (15) & 5 & 2 & \en{\ce{N^{3-}}} \\ \hline
\end{tabular}
\end{center}
\vspace{-0.25cm}
{\footnotesize * الجرمانيوم شبه فلز في المجموعة الرابعة، ويميل إلى المشاركة بالإلكترونات لا إلى تكوين أيون.}

\textbf{6-} \ans{الغازات النبيلة} (المجموعة الثامنة \en{8A}) \m{2}
\quad
\textbf{7-} العنصر الفلزي: \ans{\en{\ce{K}}} \m{1}

\textbf{(ب)} \en{$p^{+}$} = العدد الذري = \ans{15} \quad
\en{$e^{-}$} = \en{$p^{+}$} = \ans{15} \quad
\en{$n^{0}$} = العدد الكتلي $-$ العدد الذري = \en{$32 - 15$} = \ans{17} \m{6: علامتان لكل عدد}

% ============================ السؤال الرابع ============================
\question{السؤال الرابع: وضّحي المقصود بكل مما يأتي}{10 علامات: علامتان لكل تعريف}

\begin{enumerate}[label=\textbf{\arabic*)}, itemsep=0.06cm, leftmargin=*]
\item \textbf{الضغط:} القوة العمودية المؤثرة في وحدة المساحة \ (\en{$P = \dfrac{F}{A}$}).
\item \textbf{قاعدة باسكال:} الضغط المؤثر في سائل محصور ينتقل بالتساوي إلى جميع أجزاء السائل وإلى جدران الوعاء الذي يحتويه.
\item \textbf{مبدأ برنولي:} يقل ضغط المائع (السائل أو الغاز) كلما زادت سرعة جريانه.
\item \textbf{قوة الطفو:} القوة التي يؤثر بها المائع في الجسم المغمور فيه كليًّا أو جزئيًّا، واتجاهها إلى أعلى.
\item \textbf{قاعدة أرخميدس:} قوة الطفو المؤثرة في جسم مغمور في مائع تساوي وزن المائع الذي يزيحه الجسم.
\end{enumerate}

% ============================ السؤال الخامس ============================
\question{السؤال الخامس: الرافعة الهيدروليكية}{4 علامات}

حسب قاعدة باسكال: \quad
\LTRfig{$\dfrac{F_1}{A_1} = \dfrac{F_2}{A_2} \;\Rightarrow\; F_1 = \dfrac{F_2 \times A_1}{A_2} = \dfrac{24000 \times 0.4}{1.6}$} \m{2}
\par\vspace{0.1cm}
\quad القوة اللازمة: \ans{\en{$F_1 = 6000\ \mathrm{N}$}} \m{2}

% ============================ السؤال السادس ============================
\question{السؤال السادس: العلاقة بين قوة الطفو والوزن}{6 علامات: علامتان لكل حالة}

\begin{center}
\LTRfig{%
\begin{tikzpicture}[scale=0.85, >={Stealth[length=2.5mm]}]
  \foreach \x/\fb/\w/\lab/\cap in {0/1.3/0.6/{F_B > W}/{يرتفع (يطفو)}, 5/0.9/0.9/{F_B = W}/{يبقى معلقًا}, 10/0.6/0.9/{F_B < W}/{يغوص (ينزل)}} {
    \fill[blue!12] (\x,0) rectangle (\x+3.4,2.6);
    \draw[thick] (\x,3) -- (\x,0) -- (\x+3.4,0) -- (\x+3.4,3);
    \fill[orange!60] (\x+1.2,1.0) rectangle (\x+2.2,1.7);
    \draw (\x+1.2,1.0) rectangle (\x+2.2,1.7);
    \draw[->, very thick, green!50!black] (\x+1.7,1.7) -- (\x+1.7,1.7+\fb) node[right, font=\small] {$F_B$};
    \draw[->, very thick, red!70!black] (\x+1.7,1.0) -- (\x+1.7,1.0-\w) node[right, font=\small] {$W$};
    \node[font=\small] at (\x+1.7,-0.45) {$\lab$};
    \node[font=\small] at (\x+1.7,-0.95) {\foreignlanguage{arabic}{\cap}};
  }
\end{tikzpicture}}
\end{center}
\vspace{-0.15cm}
{\small \en{$F_B$}: قوة الطفو إلى أعلى، \en{$W$}: وزن الجسم إلى أسفل. إذا كانت قوة الطفو أكبر من الوزن يطفو الجسم، وإذا تساوتا يبقى معلقًا في المائع، وإذا كانت أصغر من الوزن يغوص.}

\vfill
\noindent\textbf{مع دعائي لكنّ بالنجاح} \hfill \textbf{المعلمة: دينا الشرمان}

\end{document}
